\section{Aplicaciones extensas de las CNN}

\subsection{Más allá de la clasificación: una arquitectura de propósito general}

La parte de aprendizaje de características de las CNN es altamente general: la representación de características aprendida puede usarse para múltiples tareas posteriores distintas:

\begin{figure}[H]
\centering
\includegraphics[width=0.85\textwidth]{slides-images/slide-56.jpg}
\caption{Las CNN como arquitectura de propósito general: el mismo módulo de aprendizaje de características puede sustentar clasificación, detección de objetos, segmentación semántica, control probabilístico y muchas otras tareas.\protect\footnotemark}
\end{figure}
\footnotetext{Diapositiva 56.}

\begin{itemize}
    \item \textbf{Clasificación} (Classification): determinar a qué categoría pertenece toda la imagen
    \item \textbf{Detección de objetos} (Object Detection): localizar y clasificar todos los objetos de la imagen
    \item \textbf{Segmentación semántica} (Segmentation): clasificar cada píxel
    \item \textbf{Control probabilístico} (Probabilistic Control): por ejemplo, en conducción autónoma, ir de la imagen al ángulo del volante
\end{itemize}

\subsection{Aplicación uno: detección de cáncer de mama en imágenes médicas}

Las CNN han logrado resultados revolucionarios en el análisis de imágenes médicas. Amini presentó un estudio publicado en Nature que demuestra que un sistema de IA basado en CNN \textbf{superó el desempeño de radiólogos especialistas} en la detección de cáncer de mama:

\begin{figure}[H]
\centering
\includegraphics[width=0.9\textwidth]{slides-images/slide-57.jpg}
\caption{CNN aplicadas a la detección de cáncer de mama: la curva ROC muestra que la sensibilidad y la especificidad del sistema de IA superan a las de los radiólogos. La imagen de la derecha muestra un caso de lesión mamaria detectada por la IA que un médico humano había pasado por alto.\protect\footnotemark}
\end{figure}
\footnotetext{Diapositiva 57. Referencia: McKinney y colaboradores, Nature 2020.}

\begin{importantbox}{El valor de las CNN en las imágenes médicas}
\begin{itemize}
    \item En la evaluación de la detección de cáncer de mama a 2 años y a 1 año en Estados Unidos, la curva ROC del sistema de IA superó a la del grupo de radiólogos
    \item La IA puede detectar lesiones diminutas que los médicos humanos pasan por alto, lo cual tiene un valor clínico auxiliar importante
    \item Métodos de CNN similares ya se han extendido a la detección de cáncer de piel (Esteva y colaboradores, Nature 2017) y al diagnóstico por imágenes de COVID-19
\end{itemize}
\end{importantbox}

\subsection{Aplicación dos: detección de objetos}

La \textbf{detección de objetos} (Object Detection) va más allá de la clasificación: no solo hay que identificar la categoría del objeto, sino también dar su ubicación precisa (bounding box).

\begin{figure}[H]
\centering
\includegraphics[width=0.85\textwidth]{slides-images/slide-58.jpg}
\caption{Clasificación contra detección de objetos: la clasificación solo produce una etiqueta de categoría (por ejemplo, Taxi); la detección de objetos produce a la vez la categoría y la posición del cuadro delimitador $(x, y, w, h)$.\protect\footnotemark}
\end{figure}
\footnotetext{Diapositiva 58.}

El formato de salida de la detección de objetos es un conjunto $(class, x, y, w, h)$ por cada objeto detectado:
\begin{itemize}
    \item $class$: categoría del objeto
    \item $(x, y)$: coordenadas de la esquina superior izquierda del cuadro delimitador
    \item $(w, h)$: ancho y alto del cuadro delimitador
\end{itemize}

\begin{figure}[H]
\centering
\includegraphics[width=0.85\textwidth]{slides-images/slide-59.jpg}
\caption{Detección de múltiples objetos: el modelo debe producir la categoría y el cuadro delimitador de \textbf{cada} objeto de la escena, y el número de objetos no es fijo.\protect\footnotemark}
\end{figure}
\footnotetext{Diapositiva 59.}

\begin{warningbox}{Dificultades de la detección de objetos}
La detección de objetos es mucho más difícil que la clasificación, sobre todo por lo siguiente:
\begin{enumerate}
    \item El cuadro delimitador puede aparecer en \textbf{cualquier posición} de la imagen
    \item El cuadro delimitador puede ser de \textbf{cualquier tamaño}
    \item El \textbf{número de objetos} de la escena no es fijo: puede ser cero, uno o varios
    \item Distintos objetos pueden tener distintas categorías
\end{enumerate}
Esto hace que la dimensión del espacio de salida ya no sea fija, por lo que se necesita un diseño de red especial para manejarla.
\end{warningbox}

\subsubsection{Método ingenuo: ventana deslizante}

El método más intuitivo consiste en probar en la imagen todos los cuadros delimitadores posibles, en todas las posiciones y tamaños, y clasificar con una CNN la región de imagen dentro de cada cuadro. Pero el costo de cómputo de este método es catastrófico: el número de ventanas posibles crece de manera exponencial con la posición, el tamaño y la proporción de aspecto.

\begin{figure}[H]
\centering
\includegraphics[width=0.85\textwidth]{slides-images/slide-60.jpg}
\caption{Detección de objetos ingenua: se recorren todas las ventanas posibles y cada una se envía a una CNN para clasificarla. Problema: hay demasiadas ventanas, y las combinaciones de distintas posiciones, escalas y tamaños producen una cantidad enorme de regiones candidatas.\protect\footnotemark}
\end{figure}
\footnotetext{Diapositiva 60.}

\subsubsection{La familia R-CNN}

Para resolver el problema de eficiencia, los investigadores propusieron la familia de métodos R-CNN (Region-based CNN). La idea central es generar primero un número reducido de regiones candidatas (region proposals) con métodos como la búsqueda selectiva (Selective Search), y clasificar con una CNN solo esas regiones candidatas. Los desarrollos posteriores incluyen Fast R-CNN y Faster R-CNN, que integraron aún más una red de propuesta de regiones (Region Proposal Network, RPN) dentro de la CNN, para lograr un entrenamiento de extremo a extremo.

\begin{knowledgebox}{La evolución de R-CNN a Faster R-CNN}
\begin{itemize}
    \item \textbf{R-CNN} (2014): la búsqueda selectiva genera cerca de 2,000 regiones candidatas, cada una pasa de manera independiente por una CNN para extraer características, y después se clasifica con una SVM. Es muy lento.
    \item \textbf{Fast R-CNN} (2015): primero se hace una sola extracción de características con CNN sobre toda la imagen, y luego se recortan las regiones candidatas a partir del mapa de características. Al compartir el cálculo de características, la velocidad mejora considerablemente.
    \item \textbf{Faster R-CNN} (2015): introduce una Region Proposal Network (RPN) que genera las regiones candidatas directamente sobre el mapa de características, con entrenamiento completamente de extremo a extremo.
\end{itemize}
\end{knowledgebox}

\subsection{Resumen de la sección}

El módulo de aprendizaje de características de las CNN es general y puede sustentar múltiples tareas de visión, como clasificación, detección de objetos y segmentación. En las imágenes médicas, las CNN ya alcanzan e incluso superan el nivel de los expertos humanos. La detección de objetos requiere predecir a la vez la categoría y la posición; el método ingenuo de ventana deslizante es poco eficiente, y la familia R-CNN mejoró considerablemente la eficiencia mediante el mecanismo de propuesta de regiones.

\newpage
